\subsection{06.01.2026}

\subsubsection{Ziel}

Das Ziel des Tages war es, die Probleme aus Lektion 10 zu lösen,
damit das LCD-Display zur Darstellung der Messdaten verwendet werden kann.

\subsubsection{Durchgeführte Arbeiten}

Zunächst wurde das benötigte Modul für die Ansteuerung des
LCD-Displays identifiziert und installiert.

Anschließend wurde überprüft, ob I\textsuperscript{2}C für die
Kommunikation mit dem Display aktiviert ist. Danach wurde das Skript
\texttt{lcd.py} erneut ausgeführt, um die Funktionalität des
Displays zu testen.

Zusätzlich wurde ein GitHub-Repository erstellt, um den Quellcode
für die Taupunktlüftung zentral zu verwalten. Zum Test der
GitHub-Funktionalität wurde versucht, die Skripte
\texttt{dht\_11.py} und \texttt{lcd.py} direkt vom Raspberry Pi
aus hochzuladen.

\subsubsection{Problem 1: Modul \texttt{board} nicht gefunden}

Beim Ausführen des Skripts \texttt{lcd.py} trat weiterhin die
Fehlermeldung auf, dass das Modul \texttt{board} nicht gefunden
werden konnte.

\textbf{Lösung}

Es stellte sich heraus, dass das benötigte Modul Bestandteil des
Pakets \texttt{adafruit-blinka} ist.

Nach der Installation von \texttt{adafruit-blinka} konnte das Modul
erfolgreich importiert werden.

\subsubsection{Problem 2: I\textsuperscript{2}C-Gerät nicht gefunden}

Beim erneuten Ausführen des Skripts \texttt{lcd.py} trat die
Fehlermeldung auf, dass kein I\textsuperscript{2}C-Gerät erkannt
werden konnte.

\textbf{Lösung}

Zunächst wurde überprüft, ob die richtigen Pins verwendet wurden.
Anschließend wurde kontrolliert, ob I\textsuperscript{2}C auf dem
Raspberry Pi aktiviert ist.

Nachdem I\textsuperscript{2}C aktiviert und der Raspberry Pi neu
gestartet wurde, konnte das Skript erfolgreich ausgeführt werden.

\subsubsection{Problem 3: GitHub-Authentifizierung}

Es wurde ein GitHub-Repository erstellt, um die Skripte für die
Taupunktlüftung zu verwalten und zu versionieren.

Beim Versuch, die Skripte \texttt{dht\_11.py} und
\texttt{lcd.py} vom Raspberry Pi aus hochzuladen, traten Probleme
mit der GitHub-Authentifizierung auf.

\textbf{Lösung}

Die Einrichtung der GitHub-Authentifizierung konnte an diesem Tag
noch nicht abgeschlossen werden. Bis zum nächsten Arbeitstag soll
eine geeignete Authentifizierungsmethode eingerichtet werden.

\subsubsection{Erkenntnisse}

Um I\textsuperscript{2}C-Geräte erfolgreich anzusteuern, muss
I\textsuperscript{2}C auf dem Raspberry Pi aktiviert sein.

Die Aktivierung erfolgt über folgende Schritte:

\begin{itemize}
    \item \texttt{sudo raspi-config}
    \item \texttt{Interfacing Options}
    \item \texttt{I2C}
    \item \texttt{Yes}
\end{itemize}


\subsubsection{Nächste Schritte}

Bis zur nächsten Woche soll die GitHub-Authentifizierung
eingerichtet werden, damit die Skripte erfolgreich hochgeladen
werden können.

In der nächsten Woche sollen die bestehenden Skripte weiterentwickelt
werden, um Messdaten auf dem LCD-Display darzustellen und die
Auswertung der Sensordaten vorzubereiten.
